\subsection{The two-by-three band opposite a right-zero corner}

\begin{lemma}[Proved: the top bridge group is carried by core partitions]
\label{CP2:A-two-by-three-core-partitions}
Let $M=K^1$, where
$K=\{e_{ij}:i=0,1,\ j=0,1,2\}$ and
$e_{ij}e_{k\ell}=e_{i\ell}$.  Let $N=R_3$ and let $S$
be a finite nonempty left-$N$, right-$M$ biset.
Put $R=\mathbb QM$, $T=\mathbb QN$, and $L=\mathbb Q[S]$.
The core $SK$ is a disjoint union of right-$K$ orbits $O$.
Each such orbit is a quotient of the three column labels;
write $\pi_O$ for its partition of $\{0,1,2\}$.
There is a natural left-$T$ module identification
\[
 \begin{gathered}
 \operatorname{Tor}^{R}_2(L,\mathbb Q)
   =\bigoplus_{O\subseteq SK}W_{\pi_O},\\
 W_\pi=\ker\!\left(
  \begin{aligned}
  \mathbb Q^2&\longrightarrow\mathbb Q[\{0,1,2\}/\pi],\\
  (a_1,a_2)&\longmapsto
       a_1([1]-[0])+a_2([2]-[0])
  \end{aligned}\right).
 \end{gathered}
                                                               \tag{1}
\]
In particular $\dim W_\pi=3-|O|$.  A commuting map of
sets sends an orbit to an orbit with a coarser partition,
and the corresponding map in (1) is the inclusion of
the same coordinate vectors into the new summand.
For the collage $\mathcal C_S$ there is a natural isomorphism
\[
 H_4(B\mathcal C_S;\mathbb Q)
 \cong
 \operatorname{Tor}^{T}_1\!\left(
      \mathbb Q,\operatorname{Tor}^{R}_2(L,\mathbb Q)\right).
                                                               \tag{2}
\]
\end{lemma}
\begin{proof}
For $t\in SK$ its orbit has labeled elements
$t_j=te_{ij}$ independent of $i$.  Indeed, writing
$t=se_{ab}$ gives $te_{ij}=se_{aj}$.
For every $j,\ell$ and $i$, one has
$t_je_{i\ell}=t_\ell$.  Thus two core orbits that
intersect are identical, and the only identifications
within an orbit form a partition of the three columns.
Any map commuting with $M$ preserves these labels,
which proves the stated direction of partition refinement.

Take $e=e_{00}$, $f=e_{10}$, and $q=1-e$.
The exact left-projective resolution from
Lemma~\ref{CP2:A-band-spheres} is
\[
 0\longrightarrow (Re)^2\xrightarrow{d_2}Rq
  \xrightarrow{d_1}Re\longrightarrow\mathbb Q
  \longrightarrow0,
\]
with $d_1(wq)=w(f-e)$ and
$d_{2,j}(we)=w(e_{0j}-e)$ for $j=1,2$.
After tensoring with $L$, its last three terms are
\[
 (Le)^2\longrightarrow Lq\longrightarrow Le.
                                                               \tag{3}
\]
The vector space $Le$ has one basis state $t_0$ for
each core orbit.  A basis of $Lq$ consists of
$s-se$ for all states $s$ not fixed by $e$.
This includes the non-base states of every core orbit
and all states outside $SK$.
The two columns of $d_2$ belonging to the orbit $O$
are exactly $t_1-t_0$ and $t_2-t_0$.
Their supports belong to that orbit alone.
Consequently their kernel is $W_{\pi_O}$, with no
contribution from any state outside $SK$.
This proves (1), including its naturality for every
commuting left action.

The trivial right-$T$ module has projective dimension
at most one, by the explicit right-zero resolution.
Using that resolution and (3), the double-complex
spectral sequence has
\[
 E^2_{p,q}=\operatorname{Tor}^{T}_p
       (\mathbb Q,\operatorname{Tor}^{R}_q(L,\mathbb Q)),
 \qquad 0\le p\le1,\quad 0\le q\le2.
\]
It computes the homology of the actual double-bar bridge
$E_S=B(N,S,M)$.  Its only term of total degree three
is $(p,q)=(1,2)$, and no differential enters or leaves
that term.  This gives a natural isomorphism with
$H_3(E_S;\mathbb Q)$.
The space $BM$ is a wedge of two two-spheres and $BN$
is contractible.  The natural two-block
Mayer--Vietoris sequence identifies
$H_4(B\mathcal C_S;\mathbb Q)$ with $H_3(E_S;\mathbb Q)$,
proving (2).
\end{proof}

\begin{theorem}[Proved: all states for the two-by-three band]
\label{CP2:A-two-by-three-band-right-zero-all-states}
Let $M$ be the identity-adjoined $2\times3$ rectangular
band and let $N=R_3$.  For every finite nonempty
left-$N$, right-$M$ biset, all rational degree-two
products on its collage vanish.  There is no
boundedness, Betti-number, or state-count assumption.
In particular the frozen positive order-seven source
$41$ is excluded from the eighteen-arrow budget
$(7,3,8)$.
\end{theorem}
\begin{proof}
Use the core partitions of the preceding lemma.
A singleton core orbit is an actual connecting arrow
fixed by every element of $M$.  Since $BN$ is
contractible, the dual fixed-arrow corollary
\ref{CP2:A-fixed-arrow} kills all degree-two products.
We may therefore assume that every core orbit has
two or three states.

Write $N=\{1,u,v\}$, with $u,v$ its right-zero
band elements.  The two left actors on $S$ have
common image and fix that image pointwise.  If $O$
is a two-state core orbit, $uO$ is another core
orbit with a coarser column partition.  Our exclusion
of singleton orbits forces $uO$ to have the same
two-state partition.  Every state of $uO$ is fixed
by all of $N$.  Thus each occurring two-state
partition has an actual $N$-fixed core orbit
representative $O'$.

Consider its subbridge $B(N,O',M)\to BM$.
Because $N$ acts trivially on $O'$ and $BN$ is
contractible, its homology projection is the usual
action-bar projection for $O'$.
In the resolution (3), augmentation
$\mathbb Q[O']\to\mathbb Q$ sends its degree-two
kernel to the same line $W_{\pi_{O'}}$ inside
$H_2(BM;\mathbb Q)=\mathbb Q^2$.
The three possible two-state partitions give the
three distinct lines
\[
 \mathbb Q(1,0),\qquad
 \mathbb Q(0,1),\qquad
 \mathbb Q(1,-1).
                                                               \tag{4}
\]
Hence, if at least two such partitions occur,
$H_2(E_S;\mathbb Q)\to H_2(BM;\mathbb Q)$ is
surjective.  Its cohomological pullback is injective,
so the two-block sequence makes every class in
$H^2(B\mathcal C_S;\mathbb Q)$ restrict to zero
on $BM$ and on $BN$.
Lemma~\ref{CP2:A-two-block} then kills every product.

If no two-state partition occurs, every summand in
(1) is zero, so (2) gives $H_4(B\mathcal C_S;\mathbb Q)=0$.
This also proves the desired vanishing.
It remains to consider exactly one occurring two-state
partition $\pi$.
Let $S_\pi$ be the union of all its two-state core
orbits.  It is nonempty and stable under both actions:
the right action preserves each orbit, and the left
action preserves the partition, as proved above.
Three-state core orbits have $W_\pi=0$, while every
nonzero summand in (1) belongs to $S_\pi$.
Thus the actual subcollage inclusion
\[
 B\mathcal C_{S_\pi}\longrightarrow B\mathcal C_S
                                                               \tag{5}
\]
induces an isomorphism on the left-$T$ modules in (1)
and hence an isomorphism on rational $H_4$ by (2).
It suffices to kill all products on the subcollage:
cohomological restriction along (5) is an isomorphism
in degree four and respects products.

Relabel the columns so that $\pi$ identifies columns
$1,2$ and leaves column $0$ separate.
Let $\tau$ interchange columns $1,2$, fixing both
rows and column $0$.
Every state in $S_\pi$ has equal images under those
two columns, and its right action is independent
of the row.  Therefore $\tau$ extends to an actual
involution of this subcollage, fixing $N$ and
all connecting states.

In resolution (3), $\tau$ fixes $e,f,q$.
Its semilinear lift is the identity in projective
degrees zero and one after tensoring with
$L_\pi=\mathbb Q[S_\pi]$, and interchanges the two
copies of $L_\pi e$ in degree two.
The degree-one differential is zero because the
right action is row independent.
The map
\[
 T_\pi:L_\pi e\longrightarrow L_\pi q,
       \qquad t_0\longmapsto t_1-t_0
\]
is a left-$T$ linear isomorphism: each core orbit
has exactly these two states.
Both components of $d_2$ equal $T_\pi$.
Over $\mathbb Q$ the positive eigensubcomplex of
(3) is naturally quasi-isomorphic to $L_\pi e$
in degree zero, while the negative eigensubcomplex
is $L_\pi e$ in degree two.
The positive degree-two to degree-one differential
is $2T_\pi$, so its contraction is also left-$T$ linear.

After applying $\mathbb Q\otimes_T^{\mathbf L}-$,
the right-zero projective-dimension bound shows that
$H_1(E_{S_\pi};\mathbb Q)$ is entirely positive
and $H_3(E_{S_\pi};\mathbb Q)$ entirely negative.
Thus $H_4(B\mathcal C_{S_\pi};\mathbb Q)$ is
entirely negative.
The negative eigenspace of $H_2(BM;\mathbb Q)$
is the line $\mathbb Q(1,-1)$.
The $N$-fixed two-state representative constructed
above maps onto this line from the bridge.
Consequently the dual map on that negative
cohomology eigenspace is injective.
Taking negative eigenspaces in the exact sequence
\[
 H^1(E_{S_\pi};\mathbb Q)\longrightarrow
 H^2(B\mathcal C_{S_\pi};\mathbb Q)\longrightarrow
 H^2(BM;\mathbb Q)\longrightarrow
 H^2(E_{S_\pi};\mathbb Q)
\]
shows that $H^2(B\mathcal C_{S_\pi};\mathbb Q)$
is entirely positive.
Products of two such classes are positive, whereas
all of degree four is negative.
Their products are therefore zero.
Restriction along (5) proves the same conclusion
for the original collage.

The frozen source $41$ has identity $0$ and labels
$e_{ij}=1+3i+j$.  Its forty-nine products and the
band projective resolution are independently checked in
\path{acyclic/outputs/two_row_band_row_independent_review.json}.
The present argument uses only the explicit resolution
from that check; it imposes no row-independence condition
on the original biset $S$.
\end{proof}
